\documentclass[10pt,a4paper]{extarticle}
\usepackage[a4paper,margin=10mm,top=9mm,bottom=11mm]{geometry}
% ===== Shared preamble: Class 12 Chemistry (project style, from physics-formula-sheet-source.tex) =====
\usepackage[T1]{fontenc}
\usepackage{lmodern}
\usepackage{amsmath,amssymb}
\usepackage{multicol}
\usepackage[most]{tcolorbox}
\usepackage{enumitem}
\usepackage{booktabs,tabularx,array}
\usepackage{fancyhdr}
\usepackage{xcolor}
\usepackage{microtype}
\usepackage[version=4]{mhchem}
\usepackage{tikz}
\usetikzlibrary{arrows.meta,calc}
\usepackage{pifont}
\usepackage{needspace}
\DeclareUnicodeCharacter{21CC}{\ensuremath{\rightleftharpoons}}
\DeclareUnicodeCharacter{2713}{\ding{51}}
\DeclareUnicodeCharacter{03B1}{\ensuremath{\alpha}}
\DeclareUnicodeCharacter{03B2}{\ensuremath{\beta}}
\DeclareUnicodeCharacter{2192}{\ensuremath{\to}}
\DeclareUnicodeCharacter{2610}{\ensuremath{\square}}
\DeclareUnicodeCharacter{21D2}{\ensuremath{\Rightarrow}}
\DeclareUnicodeCharacter{2194}{\ensuremath{\leftrightarrow}}
\DeclareUnicodeCharacter{2605}{\ding{72}}

% ---------- Colours: one per unit ----------
\definecolor{navy}{HTML}{1B2A4A}
\definecolor{uI}{HTML}{1F5FAD}    % 1 Solutions
\definecolor{uII}{HTML}{0E8A7E}   % 2 Electrochemistry
\definecolor{uIII}{HTML}{C2571A}  % 3 Chemical Kinetics
\definecolor{uIV}{HTML}{6A3FA0}   % 4 d- and f-Block
\definecolor{uV}{HTML}{B0204F}    % 5 Coordination Compounds
\definecolor{uVI}{HTML}{2E8B3A}   % 6 Haloalkanes & Haloarenes
\definecolor{uVII}{HTML}{8A6D00}  % 7 Alcohols, Phenols, Ethers
\definecolor{uVIII}{HTML}{2B6F8F} % 8 Aldehydes, Ketones, Acids
\definecolor{uIX}{HTML}{9A4A1E}   % 9 Amines
\definecolor{uX}{HTML}{5A5A5A}    % 10 Biomolecules
\definecolor{tipY}{HTML}{FFF4B8}
\definecolor{tipYframe}{HTML}{E0B400}
\definecolor{sqp}{HTML}{1B2A4A}
\colorlet{uc}{uI}

% ---------- Tag for questions in the official 2026-27 sample paper ----------
\newcommand{\tagbox}[2]{{\setlength{\fboxsep}{0.8pt}\colorbox{#1}{\textcolor{white}{\sffamily\bfseries\tiny\,#2\,}}}}
\newcommand{\SQP}{\tagbox{sqp}{SQP 26--27}}

% ---------- Unit banner & chapter bar ----------
\newcommand{\unitbanner}[3]{% colour, title, marks
  \colorlet{uc}{#1}%
  \par\vspace{2pt}\noindent
  \begin{tcolorbox}[enhanced,colback=#1,colframe=#1,arc=1.5pt,boxsep=1pt,left=4pt,right=4pt,top=1.5pt,bottom=1.5pt,before skip=2pt,after skip=2pt]
  \sffamily\bfseries\color{white}\footnotesize #2\hfill\scriptsize #3
  \end{tcolorbox}}
\newcommand{\chapbar}[1]{%
  \par\needspace{7\baselineskip}\noindent
  \begin{tcolorbox}[enhanced,colback=uc!18,colframe=uc,boxrule=0pt,leftrule=3pt,arc=0pt,boxsep=0.5pt,left=3pt,right=3pt,top=1pt,bottom=1pt,before skip=3pt,after skip=1.5pt]
  \sffamily\bfseries\color{uc!80!black}\footnotesize #1
  \end{tcolorbox}}
\newcommand{\sub}[1]{\par\needspace{4\baselineskip}\addvspace{2pt}\noindent{\sffamily\bfseries\scriptsize\color{uc}\strut\MakeUppercase{#1}}\par\vspace{0.5pt}}

% ---------- Boxes ----------
\newtcolorbox{kbox}{enhanced,colback=uc!5,colframe=uc,boxrule=0.6pt,arc=1.5pt,left=2pt,right=2pt,top=1pt,bottom=1pt,boxsep=1pt,before skip=2pt,after skip=2pt}
\newtcolorbox{tip}{enhanced,colback=tipY,colframe=tipYframe,boxrule=0.5pt,arc=1.5pt,left=2pt,right=2pt,top=0.5pt,bottom=0.5pt,boxsep=1pt,before skip=2pt,after skip=2pt,
  fontupper=\footnotesize}
\newtcolorbox{numbox}{enhanced,colback=white,colframe=uc!70,boxrule=0.5pt,arc=1.5pt,left=2pt,right=2pt,top=0.5pt,bottom=0.5pt,boxsep=1pt,before skip=2pt,after skip=2pt,
  borderline={0.4pt}{1.2pt}{uc!35,dashed}}
\newcommand{\Tip}{{\sffamily\bfseries\scriptsize TIP\;}}
\newcommand{\Num}{{\sffamily\bfseries\scriptsize\color{uc}STANDARD NUMERICALS\;}}

% ---------- Lists ----------
\setlist{nosep,leftmargin=8pt,labelsep=2pt}
\setlist[itemize]{label={\textcolor{uc}{\tiny$\blacktriangleright$}}}
\newenvironment{fl}{\begin{itemize}}{\end{itemize}}

% ---------- Mini graph helper (tikz) ----------
% \minigraph{width}{height}{x label}{y label}{drawing commands}
\newcommand{\minigraph}[5]{%
\begin{tikzpicture}[x=1cm,y=1cm,line cap=round,font=\tiny]
\draw[-{Stealth[length=3pt]},black!70] (0,0)--(#1,0) node[below left,inner sep=1pt]{#3};
\draw[-{Stealth[length=3pt]},black!70] (0,0)--(0,#2) node[above right,inner sep=1pt]{#4};
#5
\end{tikzpicture}}

\setlength{\columnsep}{9pt}
\setlength{\columnseprule}{0.3pt}
\def\columnseprulecolor{\color{black!25}}
\setlength{\parindent}{0pt}
\setlength{\abovedisplayskip}{1pt}\setlength{\belowdisplayskip}{1pt}
\setlength{\abovedisplayshortskip}{0pt}\setlength{\belowdisplayshortskip}{0pt}
\renewcommand{\arraystretch}{1.08}
\newcolumntype{Y}{>{\raggedright\arraybackslash}X}
\pagestyle{fancy}\fancyhf{}
\renewcommand{\headrulewidth}{0pt}

\fancyfoot[L]{\tiny\sffamily\color{black!55} Class 12 Chemistry · \sheetname · CBSE Board 2026--27 (paper 27 Feb 2027)}
\fancyfoot[R]{\tiny\sffamily\color{black!55} Page \thepage}

% ---------- Notes-specific macros ----------
\newcommand{\st}[1]{\textcolor{tipYframe!90!black}{\ifcase#1\or\ding{72}\or\ding{72}\ding{72}\or\ding{72}\ding{72}\ding{72}\fi}}
\newcommand{\ty}[1]{\tagbox{uc!75!black}{#1}}
\newcommand{\mk}[1]{{\scriptsize\color{black!55}[#1]}}
\newenvironment{rqlist}{\begin{enumerate}[leftmargin=12pt,label={\sffamily\bfseries\scriptsize\color{uc}\arabic*.},itemsep=1pt]}{\end{enumerate}}
\newcommand{\ans}{\par\hspace*{2pt}{\color{uc}$\hookrightarrow$}\;}
\newtcolorbox{skel}[1]{enhanced,colback=white,colframe=uc,boxrule=0.5pt,arc=1.5pt,left=3pt,right=3pt,top=1pt,bottom=1pt,boxsep=1pt,before skip=2pt,after skip=3pt,
  title={\sffamily\bfseries\scriptsize #1},colbacktitle=uc!85,coltitle=white,toptitle=0.5pt,bottomtitle=0.5pt}
\newcommand{\chapterhead}[4]{% colour, title, marks, subtitle
  \newpage\colorlet{uc}{#1}%
  \begin{tcolorbox}[enhanced,colback=#1,colframe=#1,arc=2pt,left=6pt,right=6pt,top=3pt,bottom=3pt,after skip=3pt]
  \color{white}\sffamily{\large\bfseries #2}\hfill{\small\bfseries #3}\\[-1pt]{\scriptsize #4}
  \end{tcolorbox}}

% orbital boxes: 0 empty, 1 up, 2 pair, 3 ligand electron pair (shaded)
\newcommand{\ob}[1]{\begin{tikzpicture}[x=2.9mm,y=2.9mm,baseline=0.3mm,line width=0.35pt]
\foreach \s [count=\i from 0] in {#1}{%
 \ifnum\s=3 \fill[uc!28] (\i,0) rectangle ++(1,1);\fi
 \draw[black!65] (\i,0) rectangle ++(1,1);
 \ifnum\s>0 \draw[-{Stealth[length=1.6pt]}] (\i+0.33,0.12)--(\i+0.33,0.88);\fi
 \ifnum\s>1 \draw[{Stealth[length=1.6pt]}-] (\i+0.67,0.12)--(\i+0.67,0.88);\fi}
\end{tikzpicture}}
\def\sheetname{Expert Q\&A}
\raggedbottom
\begin{document}
\small

% ---------- Q&A sheet macros (used after \begin{document}) ----------
\newcounter{qn}\newcounter{an}
\newcommand{\Q}{\par\addvspace{4pt}\refstepcounter{qn}\noindent\hangindent=21pt\hangafter=1\makebox[21pt][l]{\sffamily\bfseries\footnotesize\color{uc}Q\theqn.}}
\newcommand{\An}{\par\addvspace{4pt}\refstepcounter{an}\noindent\hangindent=21pt\hangafter=1\makebox[21pt][l]{\sffamily\bfseries\footnotesize\color{uc}A\thean.}}
\newcommand{\nl}{\par\noindent\hangindent=21pt\hangafter=0}
\newcommand{\secbar}[2]{\par\addvspace{4pt}\needspace{5\baselineskip}\noindent
  \begin{tcolorbox}[enhanced,colback=uc!18,colframe=uc,boxrule=0pt,leftrule=3pt,arc=0pt,boxsep=0.5pt,left=3pt,right=3pt,top=1pt,bottom=1pt,before skip=3pt,after skip=2pt]
  \sffamily\bfseries\color{uc!80!black}\footnotesize #1\hfill\scriptsize #2
  \end{tcolorbox}}
\newcommand{\qatitle}[5]{% number, title, chapters, marks, time
  \gdef\sheetname{Expert Q\&A Sheet #1 of 6: #2}
  \begin{tcolorbox}[enhanced,colback=navy,colframe=navy,arc=2pt,left=8pt,right=8pt,top=5pt,bottom=5pt]
  \color{white}\sffamily
  {\LARGE\bfseries Class 12 Chemistry — Expert Q\&A Sheet #1 of 6}\\[3pt]
  {\large #2} \hfill {\small #3}\\[2pt]
  {\footnotesize CBSE Board 2026--27 pattern · Rationalised NCERT · Paper: \textbf{Sat 27 Feb 2027} · \textbf{#4 marks · #5 min} · Questions first, full-mark answers with step marks at the end}
  \end{tcolorbox}
  \begin{tcolorbox}[enhanced,colback=tipY,colframe=tipYframe,boxrule=0.5pt,arc=1.5pt,left=4pt,right=4pt,top=1.5pt,bottom=1.5pt,before skip=3pt,after skip=3pt]
  \footnotesize\textbf{How to use.} Sit it like a test: set a timer, write on paper, no notes. Then mark yourself against the answer section, step by step. \textbf{Honest label:} every question is \textbf{PYQ-style} — the question types and data patterns that recur in CBSE papers of 2022--2026 and in NCERT — written in my own words, not copied from a specific year's paper. Step-mark splits follow the usual CBSE marking-scheme pattern and are estimates.
  \end{tcolorbox}}
\newcommand{\answerstart}{\newpage
  \begin{tcolorbox}[enhanced,colback=navy,colframe=navy,arc=2pt,left=6pt,right=6pt,top=3pt,bottom=3pt,after skip=3pt]
  \color{white}\sffamily{\large\bfseries Answers \& marking scheme}\hfill{\small written the way a full-marks script reads · {\scriptsize[\,]} = step marks}
  \end{tcolorbox}}
\newcommand{\scorebox}[1]{%
  \begin{tcolorbox}[enhanced,colback=white,colframe=navy,boxrule=0.8pt,arc=2pt,title={\sffamily\bfseries\footnotesize My score: \underline{\hspace{10mm}} / #1 \quad · \quad Time taken: \underline{\hspace{10mm}} min \quad · \quad Questions to redo / errors to log},colbacktitle=navy,coltitle=white,height fill,before skip=4pt]
  \end{tcolorbox}}

\colorlet{uc}{uIX}
\qatitle{5}{Organic Chemistry C}{Ch 9 Amines · Ch 10 Biomolecules}{42}{75}
\begin{multicols}{2}\raggedcolumns
\secbar{SECTION A · MCQ and assertion--reason}{8 × 1 = 8}
\Q The strongest base in aqueous solution is\nl (a) \ce{NH3} \quad (b) \ce{CH3NH2} \quad (c) \ce{(CH3)2NH} \quad (d) \ce{(CH3)3N} \mk{1}
\Q Which compound gives the carbylamine test?\nl (a) aniline \quad (b) N-methylaniline \quad (c) N,N-dimethylaniline \quad (d) diethylamine \mk{1}
\Q Propanamide, on treatment with bromine and aqueous NaOH, gives\nl (a) propan-1-amine \quad (b) ethanamine \quad (c) methanamine \quad (d) propanenitrile \mk{1}
\Q Which is a non-reducing sugar?\nl (a) glucose \quad (b) maltose \quad (c) lactose \quad (d) sucrose \mk{1}
\Q Pernicious anaemia is caused by deficiency of vitamin\nl (a) A \quad (b) B$_1$ \quad (c) B$_{12}$ \quad (d) K \mk{1}
\Q The base present in RNA but not in DNA is\nl (a) adenine \quad (b) thymine \quad (c) uracil \quad (d) cytosine \mk{1}
\nl\textit{For Q7--Q8 choose: (a) A and R true, R explains A; (b) both true, R does not explain A; (c) A true, R false; (d) A false, R true.}
\Q \textbf{A:} Aniline does not undergo the Friedel--Crafts reaction. \textbf{R:} Aniline forms a salt with the Lewis acid catalyst \ce{AlCl3}. \mk{1}
\Q \textbf{A:} A protein loses its biological activity on denaturation. \textbf{R:} Denaturation destroys the primary structure of the protein. \mk{1}

\secbar{SECTION B · Very short answer}{4 × 2 = 8}
\Q Write the reactions for (i) Gabriel phthalimide synthesis of ethanamine, (ii) coupling of benzenediazonium chloride with phenol. \mk{2}
\Q Arrange as directed:\nl (i) \ce{C2H5NH2}, \ce{(C2H5)2NH}, \ce{(C2H5)3N}, \ce{NH3} (increasing basic strength in water)\nl (ii) \ce{C2H5OH}, \ce{(CH3)2NH}, \ce{C2H5NH2} (increasing boiling point) \mk{2}
\Q What are reducing sugars? What happens when glucose is treated with bromine water? \mk{2}
\Q Give reasons: (i) aromatic diazonium salts are more stable than aliphatic diazonium salts; (ii) the p$K_b$ of aniline is higher than that of methylamine. \mk{2}

\secbar{SECTION C · Short answer}{4 × 3 = 12}
\Q How will you carry out these conversions?\nl (i) aniline to 4-bromoaniline \quad (ii) benzoic acid to aniline \quad (iii) aniline to iodobenzene \mk{3}
\Q Identify A, B and C:\nl \ce{CH3CH2I ->[NaCN] A ->[OH-][partial hydrolysis] B ->[NaOH + Br2] C} \mk{3}
\Q Describe the Hinsberg test for distinguishing primary, secondary and tertiary amines. \mk{3}
\Q (a) Give two differences between DNA and RNA. \mk{2}\nl (b) What type of bonding stabilises the $\alpha$-helix structure of a protein? \mk{1}

\secbar{SECTION D · Case-based}{1 × 4 = 4}
\Q \textit{Glucose, \ce{C6H12O6}, was assigned an open-chain structure from its reactions: with HI it gives n-hexane; with hydroxylamine an oxime; with bromine water gluconic acid; with acetic anhydride a pentaacetate; and with nitric acid saccharic acid. Later, several observations could not be explained by this structure, and a cyclic structure was proposed.}\\[2pt]
(i) What does the reaction with HI tell us about glucose? \mk{1}\\
(ii) Which reaction shows the presence of five \ce{-OH} groups? \mk{1}\\
(iii) State two observations that the open-chain structure cannot explain. \mk{2}\\
\textbf{OR} (iii) What are anomers? How do $\alpha$- and $\beta$-glucose differ? \mk{2}

\secbar{SECTION E · Long answer}{2 × 5 = 10}
\Q (a) Give reasons: (i) direct nitration of aniline gives a substantial amount of m-nitroaniline; (ii) aromatic primary amines cannot be prepared by Gabriel phthalimide synthesis; (iii) \ce{(CH3)2NH} is a stronger base than \ce{(CH3)3N} in aqueous solution. \mk{3}\\
(b) Give one chemical test each to distinguish (i) aniline and N-methylaniline, (ii) ethanamine and aniline. \mk{2}
\Q (a) Name the vitamins whose deficiency causes (i) scurvy, (ii) rickets. Why can't vitamin C be stored in the body? \mk{2}\\
(b) Define with one example or feature each: (i) essential amino acids, (ii) peptide linkage. \mk{2}\\
(c) Write the hydrolysis products of lactose. \mk{1}
\end{multicols}

\answerstart
\begin{multicols}{2}\raggedcolumns
\secbar{SECTION A}{1 mark each}
\An \textbf{(c) \ce{(CH3)2NH}.} Best balance of +I effect, solvation and steric factors.
\An \textbf{(a) aniline.} Only primary amines give the test.
\An \textbf{(b) ethanamine.} Hoffmann bromamide gives an amine with one carbon less.
\An \textbf{(d) sucrose.}
\An \textbf{(c) B$_{12}$.}
\An \textbf{(c) uracil.}
\An \textbf{(a).}
\An \textbf{(c).} Only the secondary and tertiary structures are destroyed; the primary structure stays intact.

\secbar{SECTION B}{2 marks each}
\An (i) Phthalimide \ce{->[alc. KOH]} potassium phthalimide \ce{->[C2H5Br]} N-ethylphthalimide \ce{->[NaOH(aq)]} \ce{C2H5NH2} + sodium phthalate \mk{1}\nl
(ii) \ce{C6H5N2+Cl- + C6H5OH ->[OH-]}\nl \ce{HO-C6H4-N=N-C6H5 + HCl} — p-hydroxyazobenzene (orange dye) \mk{1}
\An (i) \ce{NH3} $<$ \ce{C2H5NH2} $<$ \ce{(C2H5)3N} $<$ \ce{(C2H5)2NH} \mk{1}\nl
(ii) \ce{(CH3)2NH} $<$ \ce{C2H5NH2} $<$ \ce{C2H5OH} \mk{1}
\An Reducing sugars are carbohydrates that reduce Fehling's solution and Tollens' reagent; they have a free aldehydic or ketonic group (e.g.\ glucose, maltose). \mk{1}\nl
Bromine water (a mild oxidising agent) oxidises glucose to gluconic acid, showing that the carbonyl group is an aldehyde:\nl \ce{HOCH2(CHOH)4CHO ->[Br2 water] HOCH2(CHOH)4COOH} \mk{1}
\An (i) The positive charge of the arenediazonium ion is delocalised over the benzene ring by resonance; an alkyldiazonium ion has no such stabilisation. \mk{1}\nl
(ii) In aniline the lone pair on N is delocalised into the ring and is less available for protonation; in methylamine the +I effect of \ce{-CH3} increases electron density on N. Aniline is the weaker base, so its p$K_b$ is higher. \mk{1}

\secbar{SECTION C}{3 marks each}
\An (i) Aniline \ce{->[(CH3CO)2O][pyridine]} acetanilide \ce{->[Br2][CH3COOH]} 4-bromoacetanilide \ce{->[H+ or OH-][H2O]} 4-bromoaniline \mk{1}\nl
(ii) \ce{C6H5COOH ->[NH3][\Delta] C6H5CONH2 ->[Br2 + NaOH] C6H5NH2} \mk{1}\nl
(iii) \ce{C6H5NH2 ->[NaNO2 + HCl][273-278 K] C6H5N2+Cl- ->[KI] C6H5I + N2 + KCl} \mk{1}
\An \textbf{A = propanenitrile}, \ce{CH3CH2CN} \mk{1}\nl
\textbf{B = propanamide}, \ce{CH3CH2CONH2} \mk{1}\nl
\textbf{C = ethanamine}, \ce{CH3CH2NH2} (Hoffmann bromamide degradation) \mk{1}
\An The amine is shaken with benzenesulphonyl chloride (\ce{C6H5SO2Cl}, Hinsberg's reagent) in the presence of alkali.\nl
\textbf{Primary amine:} forms N-alkylbenzenesulphonamide, which is \emph{soluble in alkali} because the H on nitrogen is acidic. \mk{1}\nl
\textbf{Secondary amine:} forms N,N-dialkylbenzenesulphonamide, which has no H on nitrogen and is \emph{insoluble in alkali}. \mk{1}\nl
\textbf{Tertiary amine:} does not react. \mk{1}
\An (a) Any two \mk{1 each}: sugar is 2-deoxyribose in DNA, ribose in RNA; DNA has thymine, RNA has uracil; DNA is double-stranded (double helix), RNA is usually single-stranded; DNA carries heredity, RNA helps in protein synthesis.\nl
(b) Hydrogen bonds between the \ce{-NH} group of one amino acid residue and the $>$\ce{C=O} group of a residue in the adjacent turn of the helix. \mk{1}

\secbar{SECTION D}{4 marks}
\An (i) All six carbon atoms are joined in a straight chain. \mk{1}\nl
(ii) Acetylation with acetic anhydride to give glucose pentaacetate. \mk{1}\nl
(iii) Any two \mk{1 each}: glucose does not give the 2,4-DNP test or Schiff's test and does not form an addition product with \ce{NaHSO3}, although it has an aldehyde group; glucose pentaacetate does not react with hydroxylamine (no free \ce{-CHO}); glucose exists in two crystalline forms, $\alpha$ and $\beta$.\nl
\textbf{OR} (iii) Anomers are the two cyclic forms of a sugar that differ only in configuration at C-1, the anomeric carbon. \mk{1} In $\alpha$-glucose the \ce{-OH} at C-1 is on the right (below the ring in the Haworth structure); in $\beta$-glucose it is on the left (above the ring). \mk{1}

\secbar{SECTION E}{5 marks each}
\An (a)(i) In the strongly acidic nitrating mixture aniline is protonated to the anilinium ion, which is meta-directing. \mk{1}\nl
(ii) Aryl halides do not undergo nucleophilic substitution with the phthalimide anion (the C--X bond has partial double-bond character). \mk{1}\nl
(iii) The protonated secondary amine is better stabilised by hydrogen bonding with water (two N--H) and is less crowded than the protonated tertiary amine (one N--H); this outweighs the extra +I effect in the tertiary amine. \mk{1}\nl
(b)(i) \textbf{Carbylamine test:} on heating with \ce{CHCl3} and alcoholic KOH, aniline gives a foul-smelling isocyanide; N-methylaniline (a secondary amine) does not. \mk{1}\nl
(ii) \textbf{Azo dye test:} treat with \ce{NaNO2} + HCl at 273--278 K, then add alkaline $\beta$-naphthol. Aniline gives an orange-red dye; ethanamine gives off \ce{N2} bubbles and no dye. \mk{1}
\An (a) (i) Vitamin C; (ii) vitamin D. \mk{½ + ½} Vitamin C is water-soluble, so it is readily excreted in urine and cannot be stored. \mk{1}\nl
(b)(i) Amino acids that the body cannot synthesise and that must be supplied in the diet, e.g.\ valine, lysine. \mk{1}\nl
(ii) The amide bond \ce{-CO-NH-} formed between the \ce{-COOH} group of one amino acid and the \ce{-NH2} group of another with loss of a water molecule. \mk{1}\nl
(c) $\beta$-D-galactose and $\beta$-D-glucose. \mk{1}
\end{multicols}
\begin{tcolorbox}[enhanced,colback=tipY,colframe=tipYframe,boxrule=0.5pt,arc=1.5pt,left=4pt,right=4pt,top=2pt,bottom=2pt,before skip=4pt]
\footnotesize\textbf{Where a 70/70 script wins marks here.} (1) Basicity orders state the medium (water or gas phase) and name all three factors: +I effect, solvation, steric hindrance. (2) Diazotisation always carries the temperature (273--278 K). (3) Hinsberg and carbylamine answers say what each class of amine does. (4) Biomolecule definitions use the NCERT key words (``free aldehydic or ketonic group'', ``primary structure remains intact'', ``configuration at C-1''). (5) Conversions through acetanilide say why the \ce{-NH2} group is protected.
\end{tcolorbox}
\scorebox{42}

\end{document}
